\documentclass[10pt,a4paper]{amsart}
\usepackage[margin=19mm]{geometry}
\usepackage{amsmath,amssymb,mathtools}
\usepackage[scr=rsfs,cal=euler]{mathalfa}
\usepackage{xfrac}
\usepackage[unicode,hidelinks,bookmarks=false]{hyperref}
\newcommand{\ie}{i.e.\ }
\newcommand{\cf}{cf.\ }
\newcommand{\resp}{respectively\ }
\newcommand{\Subsection}{\subsection}
\providecommand{\nobreakdash}{\nobreak\hbox{-}}
\setlength{\emergencystretch}{2em}
\multlinegap0pt
\usepackage{tcolorbox}
\newtheorem{theorem}{Theorem}
\newcommand{\C}{\mathbb{C}}
\newcommand{\Stab}{\textnormal{Stab}}
\newcommand{\Z}{\mathbb{Z}}
\title{Explicit Analytic Continuation of Euler Products}
\author{Brandon Alberts}
\date{Source: arXiv:2406.18190v1 (CC BY 4.0)}
\begin{document}
\maketitle

\noindent\emph{Source and licence.} Explicit Analytic Continuation of Euler Products. Version arXiv:2406.18190v1 (CC BY 4.0), distributed by arXiv under the Creative Commons Attribution 4.0 International licence, http://creativecommons.org/licenses/by/4.0/, as recorded in the arXiv licence metadata for this exact version and shown on the abstract page https://arxiv.org/abs/2406.18190v1. Full licence notice: \texttt{licenses/arxiv-2406.18190-CC-BY-4.0.txt}.\par\smallskip

\noindent\textit{Editorial scope.} Counted unit: the article's theorem thm:log\_trchar (source0.tex lines 1513--1519). The article's complete original proof of it is delivered with it (source0.tex lines 1523--1627, one \texttt{proof} environment of 105 lines, which the article places away from the statement). No mathematics is written, repaired or re-derived: the statement and the proof are the article's own lines, in the article's own notation, and no retained line is substituted or re-flowed.  No further source-local context is needed: every cross-reference of every retained line resolves inside the two delivered spans.  Nothing was omitted from any retained span, so the package carries no omission record.  No retained line contains a citation, so the wrapper carries no bibliography and no bibliography entry.  No retained line contains a figure environment, an image-inclusion command or drawing code, so no image is delivered and the package declares no asset dependency.  Editorial formatting only, in addition: an AMS \texttt{amsart} preamble that restates the article's own declarations and macros used by the retained lines, namely the environments \texttt{theorem} on separate counters, together with the standard packages those lines need.  The retained environments print in this document as Theorem 1 in order of appearance, whereas the article prints the same items with the numbering of its own full document.  The complete native source of the article as distributed in the arXiv e-print for this exact version is bundled unchanged as \texttt{sources/161398/source0.tex}.
\begin{theorem}\label{thm:log_trchar}
Let $\rho:G\to GL(V)$ be a finite dimensional representation of $G$ over the $\C$-vector space $V$. Then
\[
-\log( 1 - {\rm tr}(\rho(g))z ) = -\sum_{n=1}^{\infty}\log(\det(I-\rho^{\otimes_{\text{cyc}} n}(g)z^n)),
\]
where the sum converges absolutely in the region $|z|<1$ and uniformly on the region $|z|<1-\epsilon$.
\end{theorem}

\begin{proof}[Proof of Theorem \ref{thm:log_trchar}]
We utilize a standard fact for characteristic polynomials of representations, that
\[
-\log(\det(I-\varphi(g)z)) = \sum_{k=1}^{\infty} \frac{{\rm tr}(\varphi(g^k))}{k} z^k
\]
for $|z|<1$.

Thus, it follows that
\begin{align*}
-\sum_{n=1}^{\infty}\log(\det(I-\rho^{\otimes_{\text{cyc}} n}(g)z^n)) &= \sum_{n=1}^{\infty}\sum_{k=1}^{\infty} \frac{{\rm tr}(\rho^{\otimes_{\text{cyc}} n}(g^k))}{k}z^{nk}\\
&=\sum_{k=1}^{\infty}\left(\sum_{d\mid k} \frac{{\rm tr}(\rho^{\otimes_{\text{cyc}} d}(g^{k/d}))}{k/d}\right)z^k.
\end{align*}
In order to prove the equality formally, it suffices to prove that the $k^{\rm th}$ coefficient is equal to ${\rm tr}(\rho(g))^{k}/k$.

Fix a basis for $V$, and therefore a basis of pure tensors of any tensor power labelled $\mathcal{B}_{\otimes d}$. Choose a representative of each nonzero equivalence class of pure tensor powers to form a basis for $V^{\otimes_{\text{cyc}} d}$ which we label $\mathcal{B}_{\otimes_{\text{cyc}} d}$. To compute the trace of $\rho^{\otimes_{\text{cyc}} d}(g)$, it suffices to know the coefficient of the basis element $[a]$ in the vector $\rho^{\otimes_{\text{cyc}} d}(g)[a]$, which we denote $\langle \rho^{\otimes_{\text{cyc}} d}(g)[a],[a]\rangle$, implying
\begin{align*}
{\rm tr}(\rho^{\otimes_{\text{cyc}}d}(g)) &= \sum_{\substack{[a]\in\mathcal{B}_{\otimes_{\text{cyc}} d}}} \langle\rho^{\otimes_{\text{cyc}} d}(g)[a],[a]\rangle.
\end{align*}
We note that $\rho^{\otimes_{\text{cyc}} d}(g)[a] = [\rho^{\otimes d}(g) a]$ and that $[m.a]=\zeta^m[a]$ so that the representatives of $[a]$ are all of the form $\zeta^{-m}(m.a)$. Thus
\begin{align*}
\langle \rho^{\otimes_{\text{cyc}} d}(g)[a],[a]\rangle &= \langle [\rho^{\otimes d}(g)a],[a]\rangle\\
&= \frac{1}{d}\sum_{m=0}^{d-1} \langle \rho^{\otimes d}(g)a,\zeta^{-m}(m.a)\rangle\\
&= \frac{1}{d}\sum_{m=0}^{d-1} \zeta^{-m}\langle \rho^{\otimes d}(g)a,m.a\rangle
\end{align*}
We seek to expand the summation to a sum over $\mathcal{B}_{\otimes d}$, essentially undoing the quotient defining $V^{\otimes_{\text{cyc}}d}$. A pure tensor $a=a_1\otimes \cdots \otimes a_n$ has $[a]\ne 0$ if and only if $\Stab_{\Z/d\Z}(a)=1$, which implies the orbit under the $\Z/d\Z$ action is necessarily of size $d$. We remark on two things:
\begin{enumerate}
    \item{If $a$ is a representative of the basis element $[a]$ of $V^{\otimes_{\text{cyc}} d}$ with $\Stab_{\Z/d\Z}(a) = 1$, then
    \[
    \frac{\langle \rho^{\otimes_{\text{cyc}} d}(g)[m.a],[m.a]\rangle}{\langle [m.a],[m.a]\rangle} = \frac{\langle \zeta^m\rho^{\otimes_{\text{cyc}} d}(g)[a],\zeta^m[a]\rangle}{\langle \zeta^m[a],\zeta^m[a]\rangle} = \langle \rho^{\otimes_{\text{cyc}} d}(g)[a],[a]\rangle
    \]
    is independent of the choice of representative under the $\Z/d\Z$ action. (Note that $[m.a]$ is not a basis element, so it is important that we include it in the denominator).}
    \item{If $a$ is a pure tensor with $\Stab_{\Z/d\Z}(a) = \ell\Z/d\Z$ for $\ell\ne d$, then
    \begin{align*}
    \sum_{m=0}^{d-1} \zeta^{-m}\langle \rho^{\otimes d}(g)a,m.a\rangle &= \sum_{m=0}^{\ell-1}\langle \rho^{\otimes d}(g)a,m.a\rangle \left(\sum_{j=0}^{d/\ell - 1} \zeta^{-(m+j\ell)}\right)\\
    &=\sum_{m=0}^{\ell-1}\zeta^{-m}\langle \rho^{\otimes d}(g)a,m.a\rangle \left(\sum_{j=0}^{d/\ell - 1} (\zeta^{-\ell})^j\right)\\
    &= 0,
    \end{align*}
    as we note that $\zeta^{-\ell}$ is a primitive $(d/\ell)^{\rm th}$ root of unity and $d/\ell\ne 1$.}
\end{enumerate}

Noting that the trace is the same for all bases, we partition $\mathcal{B}_{\otimes d}$ into $d$ lifts of bases for $V^{\otimes_{\text{cyc}} d}$ (which can be gotten from the orbits of a lift of $\mathcal{B}_{\otimes_{\text{cyc}}d}$) together with the collection of pure tensors $a$ with $\Stab_{\Z/d\Z}(a)\ne 1$. This produces the sum
\begin{align*}
{\rm tr}(\rho^{\otimes_{\rm cyc} d}(g))&=\frac{1}{d}\sum_{a\in\mathcal{B}_{\otimes d}}\sum_{m=0}^{d-1} \zeta^{-m} \langle\rho^{\otimes d}(g)a,m.a\rangle\\
&=\frac{1}{d}\sum_{m=0}^{d-1}\zeta^{-m} \sum_{a\in\mathcal{B}_{\otimes d}} \langle\rho^{\otimes d}(g)a,m.a\rangle.
\end{align*}
We next claim that the inner summation depends only on $\gcd(d,m)$. Indeed, for any $n$ coprime to $d$, consider the isomorphism $\{n\}:V^{\otimes d} \to V^{\otimes d}$ sending
\[
a_1\otimes \cdots \otimes a_d\mapsto a_n\otimes \cdots \otimes a_{nd}.
\]
This map certainly preserves pure tensors and commutes with $\rho^{\otimes d}(g)$ (as $\rho^{\otimes d}(g)$ acts coordinate wise). We also note that
\[
\langle \rho^{\otimes d}(g)\{n^{-1}\}a,m.\{n^{-1}\}a\rangle = \langle\{n^{-1}\} \rho^{\otimes d}(g)a,\{n^{-1}\}((nm).a)\rangle = \langle\rho^{\otimes d}(g)a,((nm).a)\rangle,
\]
where the last equality follows from $\{n^{-1}\}$ being an isomorphism which restricts to a bijection on the basis of pure tensors. Thus, we have proven
\[
\sum_{\substack{a\in\mathcal{B}_{\otimes d}\\\Stab_{\Z/d\Z}(a)=1}} \langle\rho^{\otimes d}(g)a,m.a\rangle = \sum_{\substack{a\in\mathcal{B}_{\otimes d}\\\Stab_{\Z/d\Z}(a)=1}} \langle\rho^{\otimes d}(g)a,((nm).a)\rangle
\]
whenever $\gcd(n,d)=1$, showing that this sum depends only on $\gcd(d,m)$. Partitioning as such, we find that
\begin{align*}
{\rm tr}(\rho^{\otimes_{\rm cyc} d}(g)) = \frac{1}{d}\sum_{q\mid d}\sum_{a\in\mathcal{B}_{\otimes d}} \langle\rho^{\otimes d}(g)a,q.a\rangle\left(\sum_{\substack{0\le m < d\\\gcd(m,d)=q}}\zeta^{-m} \right).
\end{align*}
The sum of roots of unity can be evaluated using the M\"obius function trick
\begin{align*}
\sum_{\substack{0\le m < d\\\gcd(m,d)=q}}\zeta^{-m} &= \sum_{\substack{0\le m < d\\q\mid m\\\gcd(m/q,d/q)=1}} \zeta^{-m}\\
&=\sum_{\substack{0\le m< d\\ q\mid m}} \zeta^{-m} \left(\sum_{\substack{\ell\mid d/q\\\ell\mid m/q}} \mu(\ell)\right)\\
&= \sum_{\ell\mid d/q} \mu(\ell) \sum_{r=0}^{d/q\ell - 1} \zeta^{-rq\ell}\\
&= \mu(d/q).
\end{align*}
Therefore we conclude that
\[
{\rm tr}(\rho^{\otimes_{\rm cyc} d}(g)) = \frac{1}{d}\sum_{q\mid d}\mu(d/q)\sum_{a\in\mathcal{B}_{\otimes d}} \langle\rho^{\otimes d}(g)a,q.a\rangle.
\]
We choose an identification $(\mathcal{B}_{\otimes q})^{d/q} \leftrightarrow \mathcal{B}_{\otimes d}$ via
\[
b_i = a_{1+iq} \otimes a_{2+iq} \otimes \cdots \otimes a_{q+iq}
\]
for each $i=1,...,d/q$. The action by $q$ on $a$ sends $b_i$ to $b_{i+1}$, taking $i$ modulo $d/q$, so we can write
\[
\langle\rho^{\otimes d}(g)a,q.a\rangle = \prod_{i=1}^{d/q} \langle \rho^{\otimes q}(g) b_i, b_{i+1}\rangle.
\]
We make use of a general fact about powers of matrices: If $A$ is an invertible matrix on a finite dimensional vector space $V$ with basis $\mathcal{B}$, then
\[
\sum_{a\in \mathcal{B}}\langle A^n a, v\rangle = \sum_{a_1,...,a_n\in \mathcal{B}} \langle A a_1,a_2\rangle\langle A a_2,a_3\rangle\cdots \langle A a_n,v\rangle.
\]
Thus,
\begin{align*}
\sum_{a\in\mathcal{B}_{\otimes d}} \langle\rho^{\otimes d}(g)a,q.a\rangle &= \sum_{b_1,...,b_{q}\in\mathcal{B}_{\otimes d/q}} \prod_{i=1}^{d/q}\langle \rho^{\otimes q}(g)b_i,b_{i+1}\rangle\\
&= \sum_{b\in\mathcal{B}_{\otimes q}}\langle \rho^{\otimes q}(g)^{d/q} b,b\rangle\\
&= {\rm tr}(\rho(g^{d/q}))^{q}
\end{align*}
and so
\[
{\rm tr}(\rho^{\otimes_{\text{cyc}} d}(g)) = \frac{1}{d}\sum_{q\mid d}\mu(d/q){\rm tr}(\rho(g^{d/q}))^q.
\]
We conclude the proof by evaluating
\begin{align*}
-\sum_{n=1}^{\infty}\log(\det(I-\rho^{\otimes_{\text{cyc}} n}(g)z^n)) &=\sum_{k=1}^{\infty}\left(\sum_{d\mid k} \frac{{\rm tr}(\rho^{\otimes_{\text{cyc}} d}(g^{k/d}))}{k/d}\right)z^k\\
&=\sum_{k=1}^{\infty}\frac{z^k}{k}\left(\sum_{d\mid k} \sum_{q\mid d}\mu(d/q) {\rm tr}(\rho((g^{k/d})^{d/q}))^q\right)\\
&=\sum_{k=1}^{\infty}\frac{z^k}{k}\left(\sum_{tyz = k}\mu(y) {\rm tr}(\rho((g^{k/ty})^{ty/t}))^t\right)\\
&=\sum_{k=1}^{\infty}\frac{z^k}{k}\left(\sum_{tyz = k}\mu(y) {\rm tr}(\rho((g^{k/t})))^t\right)\\
&=\sum_{k=1}^{\infty}\frac{z^k}{k} \sum_{t\mid k}{\rm tr}(\rho((g^{k/t})))^t\left(\sum_{y\mid k/t} \mu(y)\right)\\
&=\sum_{k=1}^{\infty}\frac{z^k}{k} {\rm tr}(\rho(g^{k/k}))^k\\
&= -\log(1 - {\rm tr}\rho(g)z).
\end{align*}
\end{proof}

\end{document}
